% Canonical workbench input module.
% No document preamble, bibliography, or machine-local path belongs here.

\section{A four-branch toroidal lift and the obstruction to a complex
\texorpdfstring{$S^8$}{S8} total space}
\label{hrftd:section}

\subsection{The exact meaning of the three-way cusp}
\label{hrftd:baseline}

The cusp in the audited threefold is controlled by the
\(\mathbb Z^2\)-periodic \(A_2\)-triangulation of \(\mathbb R^2\).  Its
quotient has one vertex, three edges, and two triangles.  The central fibre
is globally irreducible; its normalization is the degree-six del Pezzo
surface, whose six toric boundary curves are paired into three double curves.
At either of the two deepest points the smooth total threefold maps to the
disc by
\begin{equation}
 (z_0,z_1,z_2)\longmapsto t=z_0z_1z_2.
 \label{hrftd:eq:three-branch}
\end{equation}
Thus the three relevant statements are distinct: there are three local
branches at a deepest point, three global double curves, and only one global
irreducible component.  The generalized dual intersection complex is
\begin{equation}
 \mathcal T_{A_2}/\mathbb Z^2
 \cong \mathbb R^2/\mathbb Z^2=T^2.
 \label{hrftd:eq:dual-torus-two}
\end{equation}
These baseline data are transcribed from the audited canonical-source
snapshot of \(620{,}782\) bytes and \(6{,}315\) lines, with SHA--256 digest
\[
 \text{\scriptsize\ttfamily
 0696B6B7FAB83804DFEC947D45607601E5D7F57D64C6148487969FCCB6C67D2D}.
\]
line \(320\) records the periodic \(A_2\) triangulation and its
\((1,3,2)\) quotient cell counts; lines \(372\) and \(377\) record the
normalization, three double curves, and two triple points; lines
\(1431\)--\(1461\) construct the fan and its unimodular toric charts;
lines \(1647\)--\(1678\) give the quotient local models and orbit strata;
and line \(3966\) gives the deepest central-fibre equation
\(\{z_0z_1z_2=0\}\).  The line references and digest fix the exact source
edition without introducing a machine-local dependency.

The higher-dimensional construction below raises the rank of this periodic
triangulation from two to three.  It does not replace three line bundles by
four line bundles.

\subsection{The rank-three periodic fan}
\label{hrftd:periodic-fan}

Let
\[
 L=\mathbb Z^3,\qquad L_{\mathbb R}=\mathbb R^3,\qquad
 \mathbf 1=e_1+e_2+e_3.
\]
For \(m\in L\) and
\(\pi=(\pi_1,\pi_2,\pi_3)\in S_3\), put
\begin{equation}
 \sigma_{m,\pi}
 =
 \operatorname{conv}\left\{
 m,\ m+e_{\pi_1},\ m+e_{\pi_1}+e_{\pi_2},\
 m+\mathbf 1
 \right\}.
 \label{hrftd:eq:kuhn-simplex}
\end{equation}
Inside the cube \(m+[0,1]^3\), this is the region in which the coordinates,
in the order \(\pi_1,\pi_2,\pi_3\), are weakly decreasing.  The six possible
orders partition the cube, and the induced partitions agree on common
faces.  Their integral translates therefore give a locally finite
\(L\)-periodic triangulation \(\mathcal T_3\) of \(\mathbb R^3\).

Put
\[
 N'=L\oplus\mathbb Z e_4
\]
and, for every cell \(P\) of \(\mathcal T_3\), define
\begin{equation}
 \widehat P
 =
 \operatorname{cone}\{(v,1):v\in P\}
 \subset N'_{\mathbb R}.
 \label{hrftd:eq:cone-over-cell}
\end{equation}
Because all cells are bounded,
\(\widehat P\cap\widehat Q=\widehat{P\cap Q}\).  Hence the cones
\(\widehat P\), together with the zero cone, form a fan
\(\Sigma_3\).

\begin{lemma}[Unimodularity]
\label{hrftd:lemma:unimodular}
Every maximal cone of \(\Sigma_3\) is unimodular.
\end{lemma}

\begin{proof}
For the four lifted vertices of \(\sigma_{m,\pi}\), subtract the first
column \((m,1)\) from the other three columns.  Up to the orientation sign,
the determinant is
\[
 \det\left(
 e_{\pi_1},\
 e_{\pi_1}+e_{\pi_2},\
 \mathbf 1
 \right)=\pm1.
\]
The four primitive ray generators therefore form a basis of \(N'\).
Every face is unimodular as well.
\end{proof}

Let \(Y_3=X_{\Sigma_3}\) be the associated, generally non-quasi-compact,
toric fourfold.  It is smooth by
\cref{hrftd:lemma:unimodular}.  Let
\[
 M'=\operatorname{Hom}(N',\mathbb Z),\qquad
 h=(0,0,0,1)\in M',
\]
and define the toric character
\begin{equation}
 t=\chi^h:Y_3\longrightarrow\mathbb C.
 \label{hrftd:eq:toric-parameter}
\end{equation}
The character \(h\) has value one on every primitive ray generator
\((v,1)\).  Consequently
\begin{equation}
 \operatorname{div}(t)=\sum_{v\in L}E_v
 \label{hrftd:eq:reduced-central-divisor}
\end{equation}
with multiplicity one on every toric divisor \(E_v\).

\begin{proposition}[Four-branch local equation]
\label{hrftd:prop:four-branch}
If a tetrahedron \(\sigma\) has vertices \(v_0,v_1,v_2,v_3\), its toric
chart is \(\mathbb C^4\), and in its dual-character coordinates
\((z_0,z_1,z_2,z_3)\) one has
\begin{equation}
 t=z_0z_1z_2z_3.
 \label{hrftd:eq:four-branch}
\end{equation}
In particular, four distinct components of the divisor upstairs meet
transversely at the torus fixed point of \(\widehat\sigma\).
\end{proposition}

\begin{proof}
Let \(m_0,m_1,m_2,m_3\) be the basis of \(M'\) dual to
\((v_0,1),\ldots,(v_3,1)\), and set \(z_i=\chi^{m_i}\).
Unimodularity identifies the affine chart with \(\mathbb C^4\).
Since \(h\) takes value one on every ray generator,
\[
 h=m_0+m_1+m_2+m_3.
\]
Taking characters gives \eqref{hrftd:eq:four-branch}.
\end{proof}

\subsection{The analytic quotient and the abelian-threefold fibres}
\label{hrftd:quotient}

For \(\lambda\in L\), define the integral shear
\begin{equation}
 \varphi_\lambda:N'\longrightarrow N',
 \qquad
 \varphi_\lambda(y,r)=(y+r\lambda,r).
 \label{hrftd:eq:lattice-shear}
\end{equation}
It sends \(\widehat P\) to \(\widehat{P+\lambda}\), so it induces a toric
automorphism \(\Phi_\lambda\) of \(Y_3\).  These automorphisms preserve \(t\)
and satisfy
\(\Phi_\lambda\Phi_\mu=\Phi_{\lambda+\mu}\).
On the dense torus and over \(t\ne0\), their formula is
\begin{equation}
 \Phi_\lambda(x_1,x_2,x_3,t)
 =
 (t^{\lambda_1}x_1,t^{\lambda_2}x_2,t^{\lambda_3}x_3,t).
 \label{hrftd:eq:dense-action}
\end{equation}
The lattice map \eqref{hrftd:eq:lattice-shear}, rather than the expression
with possibly negative powers in \eqref{hrftd:eq:dense-action}, proves that
the action extends over the toric boundary.

Fix \(0<\epsilon<1\) and set
\[
 Y_{3,\epsilon}=\{p\in Y_3:|t(p)|<\epsilon\}.
\]
Define
\begin{equation}
 N_3=Y_{3,\epsilon}/L,\qquad
 \pi_3:N_3\longrightarrow\Delta_\epsilon.
 \label{hrftd:eq:analytic-quotient}
\end{equation}

\begin{theorem}[The rank-three toroidal cusp]
\label{hrftd:thm:toroidal-cusp}
The \(L\)-action on \(Y_{3,\epsilon}\) is free and properly discontinuous.
The quotient \(N_3\) is a smooth complex fourfold and \(\pi_3\) is proper.
For \(q\ne0\), its fibre is the abelian threefold
\[
 N_{3,q}\cong
 \mathbb C^3/(\mathbb Z^3+\tau\mathbb Z^3)
 \cong E_\tau^3,
 \qquad q=e^{2\pi i\tau},\quad \tau\in\mathbb H.
\]
\end{theorem}

\begin{proof}
A boundary point belongs to a unique torus orbit
\(O_{\widehat P}\).  If \(\Phi_\lambda\) fixes it, then
\(O_{\widehat{P+\lambda}}=O_{\widehat P}\), whence
\(P+\lambda=P\).  Since \(P\) is bounded, \(\lambda=0\).
On a fibre \(t=q\ne0\), \eqref{hrftd:eq:dense-action} is free because
\(|q|<1\) and the equalities \(q^{\lambda_i}=1\) for all \(i\) force
\(\lambda=0\).

For proper discontinuity and properness, use on the dense fibre the
logarithmic coordinate
\begin{equation}
 \operatorname{trop}_q(x)
 =
 \left(
 \frac{\log|x_1|}{\log|q|},
 \frac{\log|x_2|}{\log|q|},
 \frac{\log|x_3|}{\log|q|}
 \right).
 \label{hrftd:eq:tropical-coordinate}
\end{equation}
Equation \eqref{hrftd:eq:dense-action} translates this coordinate by
\(\lambda\).  Every orbit can therefore be moved into
\([0,1]^3\).  If
\[
 r=\sum_{i=0}^3\alpha_i v_i
\]
belongs to
\(\sigma=\operatorname{conv}(v_0,v_1,v_2,v_3)\), then
\(\alpha_i\ge0\), \(\sum_i\alpha_i=1\), and the dual-character
coordinates of \cref{hrftd:prop:four-branch} obey
\[
 \frac{\log|z_i|}{\log|q|}=\alpha_i,
 \qquad |z_i|=|q|^{\alpha_i}\le1.
\]
For \(\epsilon'<\epsilon\), let \(F_{\epsilon'}\) be the closure of the
points with \(0<|q|\le\epsilon'\) and
\(\operatorname{trop}_q(x)\in[0,1]^3\).  For each of the six tetrahedra in
that cube, the preceding identities put these points in the closed unit
polydisc \(\{|z_0|,\ldots,|z_3|\le1\}\) of its
\(\mathbb C^4\)-chart.  The finite union of those polydiscs is compact and
closed in the Hausdorff toric variety; \(F_{\epsilon'}\), being closed in
that union, is compact.  Its translates cover the dense locus, and the same
holds on \(t=0\) because every cell has an \(L\)-translate meeting the
fundamental cube.

It remains to verify proper discontinuity at the boundary, where
\(\operatorname{trop}_q\) itself is undefined.  Let \(K\subset
Y_{3,\epsilon}\) be compact.  Local compactness and the toric affine cover
give finitely many maximal-cone charts \(U_{\sigma_j}\) and compact subsets
\(K_j\subset U_{\sigma_j}\) whose interiors cover \(K\).  On the dense part
of one such compact set, write
\[
 r=\sum_{i=0}^3\alpha_iv_i,\qquad
 \alpha_i=\frac{\log|z_i|}{\log|t|},\qquad
 \sum_i\alpha_i=1.
\]
The coordinate moduli are uniformly bounded above.  When \(|t|\) is small,
this gives a uniform lower bound for every \(\alpha_i\), and the identity
\(\sum_i\alpha_i=1\) gives a uniform upper bound.  When \(|t|\) is bounded
away from zero, the product \(t=z_0z_1z_2z_3\) and the upper bounds give
uniform positive lower bounds for the \(|z_i|\), hence again uniform bounds
for every \(\alpha_i\).  Since only finitely many vertex sets occur, all
dense points of \(K\) consequently have tropical coordinates in one bounded
set \(C_K\subset\mathbb R^3\).

If \(p,\Phi_\lambda(p)\in K\) are dense, then
\(\lambda\in(C_K-C_K)\cap L\), a finite set.  If \(p\) lies on the
boundary, its orbit cell and that of \(\Phi_\lambda(p)\) lie among the
finitely many faces of the \(\sigma_j\), and they satisfy
\(P+\lambda=Q\); again only finitely many \(\lambda\) occur.  Thus a compact
set meets only finitely many of its translates, proving proper
discontinuity.  Finally, the image of \(F_{\epsilon'}\) is the entire inverse
image of \(\{|q|\le\epsilon'\}\), which proves that \(\pi_3\) is proper.
Freeness on the smooth \(Y_{3,\epsilon}\) makes \(N_3\) smooth.

For \(q\ne0\), all non-dense toric orbits lie in \(t=0\); hence
\(Y_{3,q}=(\mathbb C^*)^3\), and
\[
 N_{3,q}=(\mathbb C^*)^3/q^{\mathbb Z^3}.
\]
The map
\[
 [w_1,w_2,w_3]\longmapsto
 [e^{2\pi i w_1},e^{2\pi i w_2},e^{2\pi i w_3}]
\]
induces the asserted isomorphism from
\(\mathbb C^3/(\mathbb Z^3+\tau\mathbb Z^3)\).
\end{proof}

The preceding construction is the untwisted
\(B_0=I\), \(C(t)=0\) instance of the Tate--Mumford construction.
The general analytic construction is due to Mumford,
\emph{An analytic construction of degenerating abelian varieties over
complete rings}, Compositio Math. \textbf{24} (1972), 239--272.
No appeal to that generality is needed for the displayed quotient.

\subsection{Monodromy and the cusp inclusion on fundamental groups}
\label{hrftd:monodromy}

Use the integral basis
\[
 (a_1,a_2,a_3,b_1,b_2,b_3)
\]
of \(H_1(N_{3,q};\mathbb Z)\), where \(a_i\) is the unit period and \(b_i\)
is the \(\tau\)-period in the \(i\)-th factor.  A positive loop about
\(q=0\) sends \(\tau\) to \(\tau+1\).

\begin{proposition}[Cusp monodromy and filling map]
\label{hrftd:prop:cusp-monodromy}
The monodromy is
\begin{equation}
 M_0=
 \begin{pmatrix}I_3&I_3\\0&I_3\end{pmatrix},
 \qquad
 R=M_0-I=
 \begin{pmatrix}0&I_3\\0&0\end{pmatrix}.
 \label{hrftd:eq:cusp-monodromy}
\end{equation}
It satisfies
\[
 R^2=0,\qquad
 \operatorname{rank}R=3,\qquad
 \ker R=\operatorname{im}R
 =\mathbb Z\langle a_1,a_2,a_3\rangle.
\]
Moreover
\[
 \pi_1(N_3)\cong L\cong\mathbb Z^3,
\]
and the inclusion of a smooth fibre induces
\begin{equation}
 a_i\longmapsto0,\qquad b_i\longmapsto e_i\in L.
 \label{hrftd:eq:cusp-pi-one-map}
\end{equation}
\end{proposition}

\begin{proof}
The period change \(b_i\mapsto b_i+a_i\), \(a_i\mapsto a_i\), gives
\eqref{hrftd:eq:cusp-monodromy}; the assertions about \(R\) follow by direct
block multiplication.

The primitive rays
\[
 (0,1),\quad(e_1,1),\quad(e_2,1),\quad(e_3,1)
\]
generate \(N'\).  We record the toric fundamental-group calculation needed
here.  If \(\Sigma_0\) is a finite subfan of this smooth fan and
\(N_{\Sigma_0}\) is the sublattice generated by its primitive ray vectors,
then
\[
 \pi_1(X_{\Sigma_0})
 \cong N'/N_{\Sigma_0}.
\]
Indeed, the dense torus has fundamental group \(N'\).  Adjoining a
neighbourhood of the generic point of the divisor belonging to a primitive
ray \(v\) fills precisely the circle \(v\), so van Kampen quotients by
\(\langle v\rangle\).  The remaining orbit strata have complex codimension
at least two; general position for loops and disc homotopies in the smooth
real manifold shows that adding them changes no fundamental group.  Applying
van Kampen over all rays gives the displayed quotient.

Every loop in \(Y_3\) has compact image and is therefore contained in the
union of finitely many toric affine charts.  Enlarge their finite subfan by
the four rays displayed above.  Those rays generate \(N'\), so the finite
subfan formula makes this enlarged toric open set simply connected.  The
loop bounds there, and hence \(\pi_1(Y_3)=0\).
For \(0<u\le1\), the real one-parameter torus element
\(\vartheta_u=(1,1,1,u)\) satisfies
\[
 t(\vartheta_up)=u\,t(p).
\]
Any loop in \(Y_{3,\epsilon}\) bounds a disc in \(Y_3\).  Choosing \(u\)
small moves that entire disc into \(Y_{3,\epsilon}\), while
\(\vartheta_{u'}\), \(u\le u'\le1\), gives a homotopy back to the original
loop.  Hence \(\pi_1(Y_{3,\epsilon})=0\).  The quotient map in
\eqref{hrftd:eq:analytic-quotient} is therefore the universal cover with
deck group \(L\), so \(\pi_1(N_3)=L\).

The loop \(a_i\) is the angular loop in the \(i\)-th
\(\mathbb C^*\)-coordinate and is killed by the corresponding toric divisor.
The path representing \(b_i\) closes precisely by the deck transformation
\(\Phi_{e_i}\).  This proves \eqref{hrftd:eq:cusp-pi-one-map}.
\end{proof}

\subsection{The central fibre, its normalization, and all strata}
\label{hrftd:central-fibre}

Let
\begin{equation}
 W_3=\pi_3^{-1}(0)
 =
 \left(\bigcup_{v\in L}E_v\right)/L.
 \label{hrftd:eq:central-fibre}
\end{equation}
The toric orbit--cone correspondence identifies its open strata with
\(L\)-orbits of cells of \(\mathcal T_3\).  Its generalized dual
intersection complex is therefore
\begin{equation}
 \mathcal T_3/L\cong\mathbb R^3/L=T^3.
 \label{hrftd:eq:dual-torus-three}
\end{equation}

\begin{theorem}[Forced stratum counts]
\label{hrftd:thm:strata}
The quotient cell structure and the open torus stratification of \(W_3\)
are
\[
\begin{array}{c|c|c|c}
\text{cell}&\text{number}&\text{open stratum}&
\text{local branches}\\ \hline
\text{vertex}&1&(\mathbb C^*)^3&1\\
\text{edge}&7&(\mathbb C^*)^2&2\\
\text{triangle}&12&\mathbb C^*&3\\
\text{tetrahedron}&6&\mathrm{pt}&4.
\end{array}
\]
In particular,
\begin{equation}
 \chi(W_3)=6.
 \label{hrftd:eq:central-euler}
\end{equation}
\end{theorem}

\begin{proof}
There is one vertex orbit.  A fundamental parallelepiped of \(L\) has
normalized three-volume \(3!=6\), whereas every unimodular tetrahedron has
normalized volume one.  Hence there are six tetrahedron orbits.  Every
triangular face is incident to two tetrahedra, so
\[
 f_2=\frac{4f_3}{2}=12.
\]
Since the quotient is \(T^3\),
\[
 0=\chi(T^3)=f_0-f_1+f_2-f_3=1-f_1+12-6,
\]
and therefore \(f_1=7\).  Explicit edge representatives are
\[
 [0,e_1],\ [0,e_2],\ [0,e_3],\
 [0,e_1+e_2],\ [0,e_1+e_3],\
 [0,e_2+e_3],\ [0,\mathbf1].
\]
A \(k\)-cell gives a torus orbit of complex dimension \(3-k\) and an
intersection of \(k+1\) local branches.  Finally every
positive-dimensional \((\mathbb C^*)^j\) has Euler characteristic zero.
Additivity over the finite stratification leaves the six point strata and
proves \eqref{hrftd:eq:central-euler}.
\end{proof}

\begin{theorem}[Normalization]
\label{hrftd:thm:normalization}
\(W_3\) is globally irreducible.  The quotient map induces its normalization
\[
 \nu:E_0\longrightarrow W_3.
\]
The variety \(E_0\) is the smooth three-dimensional permutohedral toric
variety.  It has fourteen toric boundary surfaces, thirty-six invariant
curves, and twenty-four fixed points.  Under \(\nu\), these map respectively
two-to-one, three-to-one, and four-to-one onto seven double surfaces, twelve
triple curves, and six quadruple points.
\end{theorem}

\begin{proof}
The lattice \(L\) acts simply transitively on the vertices and therefore
transitively on the components \(E_v\).  The quotient \(W_3\) is irreducible.
Every cell has a vertex and can be translated to one containing zero, so the
map \(E_0\to W_3\) is surjective.  Only finitely many cells contain zero,
hence the map is finite, and it is injective on the dense torus.  The star
fan at zero is complete and unimodular, so \(E_0\) is smooth, compact, and
normal.  The finite birational map is consequently the normalization.

The rays of the star fan are
\[
 \pm\sum_{i\in I}e_i,
 \qquad
 \varnothing\ne I\subseteq\{1,2,3\}.
\]
These are the fourteen rays of the \(A_3\) braid fan.  Its maximal cones are
indexed by the twenty-four permutations of four symbols; hence the fan is the
permutohedral fan.  There are \(24\) maximal cones.  Each has three
two-dimensional faces and each such face is shared by two maximal cones, so
there are \(24\cdot3/2=36\) two-dimensional cones.  Together with the
fourteen rays this gives the asserted orbit counts.

An edge has two vertices, a triangle has three, and a tetrahedron has four.
The normalization preimage of the corresponding generic quotient stratum
therefore has respectively two, three, or four points.  Division of the
normalization counts by these degrees gives
\[
 14/2=7,\qquad 36/3=12,\qquad 24/4=6.
\]
\end{proof}

\begin{remark}[Local branches are not global components]
\label{hrftd:rem:local-not-global}
At each quadruple point the four local branches are four germs of the same
global irreducible central fibre.  The conductor in the normalization is the
union of fourteen toric boundary surfaces, paired into seven quotient
surfaces.  No canonical four-term decomposition of a vector bundle or line
bundle follows from this local branch number.
\end{remark}

More generally, let \(\bar L\cong\mathbb Z^3\) and let
\(B_0:\bar L\hookrightarrow L\) have finite index \(d\).  Replacing
\(\lambda\) by \(B_0\lambda\) in \eqref{hrftd:eq:lattice-shear} gives
\(d\) vertex orbits and hence \(d\) global central components.  The stratum
counts are
\begin{equation}
 (f_0,f_1,f_2,f_3)=d(1,7,12,6).
 \label{hrftd:eq:index-d-counts}
\end{equation}

\begin{corollary}[Euler characteristic as the determinant of the shear]
\label{hrftd:cor:euler-shear-determinant}
Choose oriented integral bases of \(\bar L\) and \(L\), and let
\([B_0]\in\operatorname{Mat}_{3\times3}(\mathbb Z)\) be the matrix of the
injective shear.  The associated central fibre satisfies
\begin{equation}
 \chi(W_{3,B_0})
 =3!\,[L:B_0\bar L]
 =3!\,|\det[B_0]|.
 \label{hrftd:eq:euler-shear-determinant}
\end{equation}
The right side is unchanged by independent unimodular changes of the two
marked bases.
\end{corollary}

\begin{proof}
The index of the image of a full-rank integral matrix is
\([L:B_0\bar L]=|\det[B_0]|\).  By
\eqref{hrftd:eq:index-d-counts}, the quotient has \(6d\) tetrahedron orbits.
Their zero-dimensional torus strata are the only strata with nonzero Euler
characteristic, so \(\chi(W_{3,B_0})=6d=3!d\).  Replacing the source and
target bases multiplies the determinant by units in \(\mathbb Z\), and hence
does not change its absolute value.
\end{proof}

If a polarization additionally identifies \([B_0]\) with a symmetric
positive-definite matrix, then its determinant is the cubic Jordan norm on
that real symmetric subcone of
\(\operatorname{Herm}_3(\mathbb C)_{\geq0}\).  Without those additional
hypotheses \([B_0]\) is only an integral lattice map, not a positive
Hermitian matrix.  In either case the scalar in
\eqref{hrftd:eq:euler-shear-determinant} does not determine the fan, the
period correction, the singular filling, or the total space.

Thus four global components occur only after the independent choice \(d=4\);
four local branches already occur for \(d=1\).  A nontrivial Fable-style
twist would additionally require holomorphic translations
\(c_\lambda(t)\) satisfying the exact group cocycle, extension, freeness,
and properness identities.

\subsection{Two independent obstructions to a complex
\texorpdfstring{$S^8$}{S8} globalization}
\label{hrftd:s8-obstructions}

\begin{theorem}[Euler obstruction to the literal three-special-fibre pattern]
\label{hrftd:thm:euler-obstruction}
Let \(g:M^8\to S^2\) be a proper triangulable map which is a locally trivial
\(T^6\)-bundle away from finitely many critical values.  If one singular
fibre is \(W_3\) and every other singular fibre is a finite free quotient of
\(T^6\), then
\[
 \chi(M)=6.
\]
In particular, \(M\) is not homeomorphic to \(S^8\).
\end{theorem}

\begin{proof}
Choose a triangulation for \(g\) in which every critical value is a vertex.
Additivity of compactly supported Euler characteristic over the open
regular stratum and the critical singleton strata, proved cell by cell in
this finite triangulation, gives
\[
 \chi(M)
 =
 \chi_c(S^2\setminus\Sigma)\chi(T^6)
 +\sum_{p\in\Sigma}\chi(g^{-1}(p)).
\]
Here multiplicativity on the regular stratum follows by subdividing its
locally trivial bundle over base cells.  Since \(M\) is compact,
\(\chi_c(M)=\chi(M)\), and since \(\chi(T^6)=0\), the first term vanishes.
By \eqref{hrftd:eq:central-euler}, the cusp contributes \(6\).
A finite free quotient by a group \(G\) has Euler characteristic
\(\chi(T^6)/|G|=0\).  Therefore
\(\chi(M)=6\), whereas \(\chi(S^8)=2\).
\end{proof}

Consequently, retaining the cusp \(W_3\) in a real \(S^8\) candidate would
require additional genuinely singular fibres whose Euler characteristics sum
to \(-4\).  Smooth multiple torus fibres cannot supply them.

\begin{theorem}[No smooth complex \(S^8\) total space]
\label{hrftd:thm:no-almost-complex-s8}
A smooth manifold homotopy equivalent to \(S^8\) admits no almost-complex
structure.  Hence a complex structure on the regular locus of a putative
\(T^6\)-fibration on \(S^8\) cannot extend through all its singular fillings.
\end{theorem}

\begin{proof}
Suppose \(E=(TM,J)\) is a complex rank-four bundle whose underlying real
bundle is \(TM\).  Orient \(M\) by \(J\), and let
\(u\in H^8(M;\mathbb Z)\) evaluate to one.  Since
\[
 H^2(M;\mathbb Z)=H^4(M;\mathbb Z)=H^6(M;\mathbb Z)=0,
\]
one has
\[
 c_1(E)=c_2(E)=c_3(E)=0,\qquad
 c_4(E)=e(TM)=\chi(M)u=2u.
\]
For the underlying real bundle,
\begin{equation}
 p_2(TM)
 =c_2(E)^2-2c_1(E)c_3(E)+2c_4(E)
 =4u.
 \label{hrftd:eq:p-two-contradiction}
\end{equation}
On the other hand \(H^4(M)=0\) gives \(p_1(TM)=0\), and the middle
cohomology vanishes, so \(\operatorname{sign}(M)=0\).  Hirzebruch's
eight-dimensional signature formula gives
\[
 0
 =
 \left\langle
 \frac{7p_2(TM)-p_1(TM)^2}{45},[M]
 \right\rangle,
\]
which forces \(p_2(TM)=0\).  This contradicts
\eqref{hrftd:eq:p-two-contradiction}.  For the standard sphere the same
vanishing follows immediately from
\(TS^8\oplus\varepsilon^1_{\mathbb R}\cong\varepsilon^9_{\mathbb R}\).
\end{proof}

If every local filling were a smooth complex toroidal fourfold and every
overlap map were biholomorphic, these local complex tangent bundles would glue
to an almost-complex structure on all of \(S^8\), contradicting
\cref{hrftd:thm:no-almost-complex-s8}.  At least one filling, transition, or
total-space singularity must therefore leave the smooth complex category.

\subsection{An explicit integral monodromy candidate and its failure}
\label{hrftd:finite-monodromy}

Put
\[
 S=
 \begin{pmatrix}0&-1\\1&0\end{pmatrix},
 \qquad
 T=
 \begin{pmatrix}1&1\\0&1\end{pmatrix},
 \qquad
 Q=S^{-1}T^{-1}
 =
 \begin{pmatrix}0&1\\-1&1\end{pmatrix}.
\]
Then
\begin{equation}
 S^4=I,\qquad Q^6=I,\qquad SQ=T^{-1}.
 \label{hrftd:eq:two-by-two-relations}
\end{equation}
On
\(\Lambda=\mathbb Z^3\otimes\mathbb Z^2\cong\mathbb Z^6\), define
\begin{equation}
 M_1=I_3\otimes S,\qquad
 M_2=I_3\otimes Q,\qquad
 M_0=I_3\otimes T.
 \label{hrftd:eq:finite-monodromies}
\end{equation}
They are integral and symplectic, and
\[
 M_1M_2M_0=I.
\]
After reordering the tensor-product basis as
\((a_1,a_2,a_3,b_1,b_2,b_3)\), \(M_0\) is exactly the cusp matrix
\eqref{hrftd:eq:cusp-monodromy}.  Thus
\eqref{hrftd:eq:finite-monodromies} is an explicit local system over a
thrice-punctured sphere with finite orders \(4,6\) and a rank-three
unipotent cusp.

\begin{proposition}[Fixed-point obstruction]
\label{hrftd:prop:finite-fixed-points}
No translation makes either finite linear action \(M_1\) or \(M_2\) free on
\(T^6\).
\end{proposition}

\begin{proof}
Neither \(S\) nor \(Q\) has eigenvalue one.  More explicitly,
\[
 \det(I-S)=2,\qquad \det(I-Q)=1.
\]
Consequently \(I-M_j:T^6\to T^6\) is a surjective isogeny for \(j=1,2\).
For any translation parameter \(c\in T^6\), the fixed-point equation for
\[
 x\longmapsto M_jx+c
\]
is
\[
 (I-M_j)x=c,
\]
and it always has a solution.  Therefore every affine lift has a fixed point.
\end{proof}

This is an integral symplectic homology local system compatible with a
diagonal rank-three period ansatz.  The three matrices alone do not define a
polarized variation of Hodge structure or a holomorphic family
\(E_\tau^3\); those require an equivariant period map in addition.  Any
isotrivial affine quotient filling built from either displayed finite
generator is singular by the fixed-point calculation.  A successful
alternative would need finite
monodromies with invariant torus directions that admit fixed-point-free
affine lifts, while retaining the product relation and the rank-three cusp.
It would still have to overcome both obstructions in
\cref{hrftd:s8-obstructions}.

\subsection{Exact replication requirements for any real singular
\texorpdfstring{$T^6\to S^8\to S^2$}{T6 to S8 to S2} candidate}
\label{hrftd:replication}

The preceding results exclude the direct smooth holomorphic construction, not
an arbitrary real stratified map.  Let
\[
 \Sigma=\{p_0,\ldots,p_s\}\subset S^2,\qquad
 B^\circ=S^2\setminus\Sigma.
\]
A candidate must first specify an affine integral monodromy
\begin{equation}
 \widetilde\rho:\pi_1(B^\circ)\longrightarrow
 \operatorname{GL}(6,\mathbb Z)\ltimes T^6
 \label{hrftd:eq:affine-monodromy}
\end{equation}
with based peripheral matrices satisfying
\[
 M_0M_1\cdots M_s=I.
\]
The corresponding affine peripheral transformations
\(\widetilde M_i=(M_i,c_i)\) must satisfy the stronger identity
\begin{equation}
 \widetilde M_0\widetilde M_1\cdots\widetilde M_s=(I,0),
 \qquad
 c_0+M_0c_1+\cdots+M_0\cdots M_{s-1}c_s=0
 \quad\text{in }T^6.
 \label{hrftd:eq:affine-product}
\end{equation}
Choose pairwise disjoint small closed discs \(D_i\) around the \(p_i\), and
put
\[
 B=S^2\setminus\bigcup_i\operatorname{int}D_i,\qquad
 J=g^{-1}(B).
\]
It must then give each compact filling \(N_i=g^{-1}(D_i)\), its boundary
mapping torus
\[
 A_i=J\cap N_i,
\]
and the actual based inclusion maps
\[
 A_i\longrightarrow J,\qquad A_i\longrightarrow N_i.
\]
Here one fixes a base point in \(J\), a base point in each \(A_i\), and a
whisker in \(J\) joining them; changing a whisker only conjugates the
corresponding homomorphism and does not change the pushout below.

\begin{proposition}[Fundamental-group replication map]
\label{hrftd:prop:van-kampen}
The fundamental group of the glued total space is the iterated pushout
\begin{equation}
 \pi_1(M)
 \cong
 \pi_1(J)
 *_{\pi_1(A_0)}\pi_1(N_0)
 *_{\pi_1(A_1)}\cdots
 *_{\pi_1(A_s)}\pi_1(N_s),
 \label{hrftd:eq:van-kampen}
\end{equation}
where both homomorphisms out of each \(\pi_1(A_i)\) are the literal based
inclusion homomorphisms.  If, in addition, every
\(\pi_1(A_i)\to\pi_1(N_i)\) is surjective, then
\begin{equation}
 \pi_1(M)
 \cong
 \pi_1(J)\Big/
 \Big\langle\!\Big\langle
 (j_{J,i})_*\ker\{\pi_1(A_i)\to\pi_1(N_i)\}:0\le i\le s
 \Big\rangle\!\Big\rangle .
 \label{hrftd:eq:van-kampen-surjective}
\end{equation}
Consequently a proposed \(S^8\) construction must either compute the full
pushout in \eqref{hrftd:eq:van-kampen}, or prove the stated surjectivities
and prove that the displayed normal generators kill \(\pi_1(J)\).
\end{proposition}

\begin{proof}
Apply van Kampen successively to
\[
 M=J\cup N_0\cup\cdots\cup N_s
\]
with connected collar intersections \(A_i\).  Each attachment quotients by
the relations identifying the two images of every element of
\(\pi_1(A_i)\), which is precisely the iterated amalgamated product
\eqref{hrftd:eq:van-kampen}.  If the map to \(\pi_1(N_i)\) is surjective,
then every generator coming from \(\pi_1(N_i)\) can be represented by an
element of \(\pi_1(A_i)\).  Eliminating those generators leaves exactly the
normal closure of the image in \(\pi_1(J)\) of the kernel of the map to
\(\pi_1(N_i)\).  Iterating proves
\eqref{hrftd:eq:van-kampen-surjective}.
\end{proof}

For homology, define the literal Mayer--Vietoris map
\begin{equation}
 \alpha_q
 =
 \bigoplus_i\bigl((j_{J,i})_*,-(j_{N_i})_*\bigr):
 \bigoplus_iH_q(A_i;\mathbb Z)
 \longrightarrow
 H_q(J;\mathbb Z)\oplus\bigoplus_iH_q(N_i;\mathbb Z).
 \label{hrftd:eq:mv-map}
\end{equation}

\begin{theorem}[Integral homology-sphere criterion]
\label{hrftd:thm:mv-criterion}
Assume that \(J\) and the \(N_i\) are compact, connected,
collar-compatible eight-manifolds with nonempty boundary, so their top
homology vanishes.  The glued space has the integral homology of \(S^8\) if
and only if
\[
\begin{array}{c|c}
q&\text{condition on }\alpha_q\\ \hline
0&\alpha_0\text{ is injective and }
   \operatorname{coker}\alpha_0\cong\mathbb Z,\\
1\le q\le6&\alpha_q\text{ is an isomorphism},\\
7&\alpha_7\text{ is surjective and }
   \ker\alpha_7\cong\mathbb Z.
\end{array}
\]
\end{theorem}

\begin{proof}
Insert \eqref{hrftd:eq:mv-map} into the Mayer--Vietoris long exact sequence.
For every \(1\le q\le7\), exactness gives the short exact sequence
\[
 0\longrightarrow\operatorname{coker}\alpha_q
 \longrightarrow H_q(M;\mathbb Z)
 \longrightarrow\ker\alpha_{q-1}
 \longrightarrow0.
\]
Thus simultaneous vanishing in degrees \(1,\ldots,7\) is equivalent to
\(\alpha_0\) being injective, \(\alpha_q\) being an isomorphism for
\(1\le q\le6\), and \(\alpha_7\) being surjective.  In degree zero,
\(H_0(M)\cong\operatorname{coker}\alpha_0\), which gives the remaining
degree-zero condition.  Since the pieces have zero \(H_8\), exactness at
\[
 0\longrightarrow H_8(M)\longrightarrow
 \bigoplus_iH_7(A_i)\xrightarrow{\alpha_7}
 H_7(J)\oplus\bigoplus_iH_7(N_i)\longrightarrow H_7(M)
 \longrightarrow0
\]
identifies \(H_8(M)\) with \(\ker\alpha_7\).  Hence
\(H_8(M)\cong\mathbb Z\) is exactly the last kernel condition.
\end{proof}

If \(A_i\) is the mapping torus of a torus automorphism with integral
monodromy \(M_i\), its integral homology is fixed by the Wang exact sequence:
\begin{equation}
 0\longrightarrow
 \operatorname{coker}\left(I-\bigwedge^qM_i\right)
 \longrightarrow H_q(A_i;\mathbb Z)
 \longrightarrow
 \ker\left(I-\bigwedge^{q-1}M_i\right)
 \longrightarrow0.
 \label{hrftd:eq:wang}
\end{equation}
The right-hand group is free, so this sequence splits as a sequence of
abelian groups, though not canonically.  Equations
\eqref{hrftd:eq:affine-monodromy},
\eqref{hrftd:eq:van-kampen}, \eqref{hrftd:eq:mv-map}, and
\eqref{hrftd:eq:wang}, together with the actual tangent clutching maps, are
the minimum replication data.  For an \(S^8\) total space those clutching
maps must give
\[
 p_1=p_2=0,\qquad e=2u,
\]
while preventing the regular-locus complex structure from extending, as
required by \cref{hrftd:thm:no-almost-complex-s8}.

\subsection{The Hermitian determinant boundary is not the four-branch fibre}
\label{hrftd:hermitian}

Let
\[
 \mathcal C_3=\operatorname{Herm}_3(\mathbb C)_{\ge0}.
\]
The real dimension of the positive-semidefinite rank-\(r\) stratum is
\[
 2\cdot3r-r^2.
\]
Indeed, write \(H=BB^*\) with \(B\) a full-rank complex
\(3\times r\) matrix; the entire redundancy is the right action of
\(U(r)\).  Thus the rank \(3,2,1,0\) strata have real dimensions
\begin{equation}
 9,\quad8,\quad5,\quad0.
 \label{hrftd:eq:rank-dimensions}
\end{equation}
The rank-two stratum is smooth because the derivative of the determinant is
the trace pairing with \(\operatorname{adj}(H)\), which is nonzero exactly
in rank two.  The singular locus is the rank-at-most-one locus.

\begin{proposition}[The real Hermitian boundary compactifies to \(S^8\)]
\label{hrftd:prop:hermitian-compactification}
The determinant-zero boundary
\[
 \partial\mathcal C_3
 =\{H\in\operatorname{Herm}_3(\mathbb C):H\ge0,\ \det H=0\}
\]
is homeomorphic to \(\mathbb R^8\).  Its one-point compactification is
therefore \(S^8\), with the rank strata in
\eqref{hrftd:eq:rank-dimensions} giving a non-smooth stratification.
\end{proposition}

\begin{proof}
Intersect \(\mathcal C_3\) with the affine hyperplane \(\operatorname{tr}H=1\):
\[
 K=\{H\ge0:\operatorname{tr}H=1\}.
\]
This is a compact convex body with nonempty interior in the
eight-dimensional affine space
\(\{H=H^*:\operatorname{tr}H=1\}\).  Radial projection from the interior
point \(I_3/3\) gives a homeomorphism from \(K\) to the closed
eight-ball.  Explicitly, after choosing a Euclidean norm on the trace-zero
translation space, every ray
\[
 I_3/3+sU,\qquad \|U\|=1,\quad s\ge0,
\]
meets \(\partial K\) once, at a radius \(r(U)>0\); compactness and convexity
make \(r\) continuous, and
\((\rho,U)\mapsto I_3/3+\rho r(U)U\), \(0\le\rho\le1\), is the required
radial homeomorphism after collapsing \(\{0\}\times S^7\) to the centre.
Thus \(\partial K\cong S^7\).

Every nonzero \(H\in\partial\mathcal C_3\) has the unique expression
\[
 H=rP,\qquad r=\operatorname{tr}H>0,\qquad P\in\partial K.
\]
Adding \(H=0\) identifies
\[
 \partial\mathcal C_3
 \cong
 \bigl([0,\infty)\times S^7\bigr)/(\{0\}\times S^7)
 \cong\mathbb R^8.
\]
The Alexandroff one-point compactification of \(\mathbb R^8\) is \(S^8\).
\end{proof}

\begin{theorem}[Non-SNC determinant boundary]
\label{hrftd:thm:hermitian-not-snc}
At a rank-one point, the semialgebraic boundary of \(\mathcal C_3\) has a
quadratic-cone normal slice rather than a union of smooth crossing faces.
The ambient real determinant hypersurface, and its complexification, are not
normal-crossings hypersurfaces there.  Neither that rank-one germ nor the
germ at the zero matrix is the four-branch germ in
\eqref{hrftd:eq:four-branch}.
\end{theorem}

\begin{proof}
Unitary conjugation and positive rescaling reduce any rank-one point to
\[
 H_0=\operatorname{diag}(1,0,0).
\]
Near \(H_0\), write
\[
 H=
 \begin{pmatrix}a&b^*\\b&C\end{pmatrix},
 \qquad
 a\in\mathbb R,\quad b\in\mathbb C^2,\quad
 C\in\operatorname{Herm}_2(\mathbb C).
\]
On \(a>0\), the real-analytic coordinate change
\[
 (a,b,C)\longmapsto(a,b,S),\qquad
 S=C-a^{-1}bb^*
\]
is a local diffeomorphism.  The Schur-complement identities are exact:
\begin{equation}
 \det H=a\det S,\qquad
 H\ge0\quad\Longleftrightarrow\quad S\ge0.
 \label{hrftd:eq:schur}
\end{equation}
Writing
\[
 S=\begin{pmatrix}u&x+iy\\x-iy&v\end{pmatrix}
\]
gives the exact semialgebraic normal boundary
\begin{equation}
 uv-x^2-y^2=0,\qquad u\ge0,\quad v\ge0.
 \label{hrftd:eq:quadratic-cone}
\end{equation}
The defining quadratic form has rank four.  A product of two real linear
forms has quadratic rank at most two, so the ambient real-algebraic
hypersurface is not a union of smooth real hyperplane branches; imposing the
displayed positive-semidefinite inequalities cannot turn its quadratic cone
into transverse smooth faces.

After complexification, the normal slice is
\[
 s_{11}s_{22}-s_{12}s_{21}=0.
\]
All four first derivatives vanish only at the origin.  A product of two
linear forms would be singular along their codimension-two intersection, so
the quadratic is irreducible and is not a coordinate crossing.  At the zero
matrix the \(3\times3\) determinant is homogeneous of degree three and
therefore has multiplicity three.  The germ
\(z_0z_1z_2z_3=0\) has multiplicity four.  Neither germ is the other.
\end{proof}

There is also a dimension obstruction.  The toroidal central fibre \(W_3\)
has complex dimension three and real dimension six, while the dense rank-two
boundary of \(\mathcal C_3\) has real dimension eight.  The one-point
compactification in \cref{hrftd:prop:hermitian-compactification} is a
stratified \(S^8\), but it is not the central divisor of the complex
fourfold \(N_3\).

The two occurrences of the number three are independent:
\(\det_{\operatorname{Herm}_3}\) has three eigenvalue factors, whereas a
semistable degeneration of complex three-tori has \(3+1=4\) local branches
at a deepest point because one includes the base parameter.  Moving to
\(\operatorname{Herm}_4\) makes the determinant quartic but changes the real
ambient dimension to \(16\) and the natural compactified boundary dimension
to \(15\); it does not produce an \(S^8\).

\subsection{Exact conclusion and permissible carriers}
\label{hrftd:conclusion}

The local higher-rung toroidal object is exact:
\[
 \boxed{
 \text{rank-three periodic fan}
 \quad\Longrightarrow\quad
 t=z_0z_1z_2z_3,\quad
 \mathcal D(W_3)\cong T^3,\quad
 (f_0,f_1,f_2,f_3)=(1,7,12,6).
 }
\]
It lives on the smooth noncompact complex fourfold \(N_3\to\Delta\), not on
a complex \(S^8\).  A compact complex globalization must have different
topology.  A real \(S^8\) may retain the complex-torus structure only on its
regular locus and requires explicitly non-complex or singular fillings
satisfying every map in \cref{hrftd:replication}.  The Hermitian Gram cone
provides a separate real stratified \(S^8\) boundary, not an SNC toroidal
fibre.  Stable complex data are also permissible because
\[
 TS^8\oplus\varepsilon^2_{\mathbb R}
 \cong\varepsilon^{10}_{\mathbb R},
\]
but this stable structure cannot destabilize to an almost-complex structure
on \(TS^8\).

Finally, the six tetrahedron orbits are naturally indexed by the six
permutations in \(S_3\); neither this integer six nor the rank-six lattice
\(H_1(T^6;\mathbb Z)\) is by itself a cyclic group
\(\mathbb Z/6\mathbb Z\).  Any cyclic-order-six or gauge-centre
interpretation requires an additional exact homomorphism.
